% buffer-shell.tex — Chapter on the single-plaquette buffer-shell bound.
%
% Linked code / Lean:
%   - computations/buffer_shell.py
%   - lean/YMGap/BufferShell.lean

\chapter{The Single-Plaquette Buffer-Shell Bound}
\label{ch:buffer-shell}

\noindent\emph{Two named bounds on the heat-kernel mass in the
buffer annulus: the spec-shape bound with prefactor $g^3$ and the
uniform-in-$g$ bound used downstream.}

This chapter establishes the single-plaquette buffer-shell estimate
that drives the per-plaquette content of the overlap-measure
analysis in \Cref{ch:overlap-measure}.  The setting is one
plaquette $p$ on the scale-$k$ lattice, with link variables
distributed according to the Wilson-action Gibbs measure at coupling
$g_k$; the buffer annulus is the configuration set where the
plaquette holonomy $U_p$ sits at geodesic distance from the identity
in the indeterminate annulus between the small-field and large-field
regions.

\section{Setup and the buffer annulus}
\label{sec:buffer-setup}

Throughout this chapter we work in the canonical-$\kappa$ convention
of \Cref{ch:heat-kernel}, Convention~\ref{conv:canonical-kappa},
unless otherwise stated.  The plaquette holonomy is parameterized
by its eigenphase $\theta \in [0, 2\pi)$ through the maximal-torus
representative
$U_p = \mathrm{diag}(e^{i\theta/2}, e^{-i\theta/2})$
(\Cref{ch:su2-rep-theory}, equation~\eqref{eq:torus-element}); the
geodesic distance from the identity is $d(U_p, 1) = \theta$ when
restricted to the principal branch $\theta \in [0, \pi]$.

\marginkey{The buffer annulus}%

\begin{definition}[Buffer annulus]
\label{def:buffer-annulus}
Let $\delta_k > 0$ be the small-field smallness parameter at scale
$k$.  The \emph{buffer annulus} at scale $k$ is
\begin{equation}
  \label{eq:buffer-annulus}
  A(\delta_k) \;:=\; \bigl\{ U \in \SUtwo : \tfrac{\delta_k}{2} < d(U, 1) < 2\,\delta_k \bigr\}
  \;\subseteq\; \SUtwo.
\end{equation}
\end{definition}

The buffer-shell \emph{measure} at scale $k$ is the probability of
$U_p \in A(\delta_k)$ under the heat-kernel Gibbs measure $\muk$:
\begin{equation}
  \label{eq:buffer-shell-prob}
  P_{\mathrm{shell}}(\delta_k, g_k)
  \;:=\;
  \frac{1}{\pi}
  \int_{\delta_k / 2}^{2\,\delta_k}
    K_t(\theta)\, \sin^2(\theta/2)\, d\theta ,
\end{equation}
with $t = t(g_k)$ fixed by the chosen convention and $K_t$ the
$\SUtwo$ heat kernel of \Cref{ch:heat-kernel}, equation~\eqref{eq:heat-kernel-su2}.
Equation~\eqref{eq:buffer-shell-prob} is the Weyl-integration
realization of $\muk(A(\delta_k))$ on a single plaquette.

\section{Two named bounds and their shapes}
\label{sec:buffer-bounds}

The buffer-shell probability admits two named upper bounds, each
useful in a different regime.

\marginkey{Spec-shape bound}%

\begin{theorem}[Spec-shape bound]
\label{thm:buffer-spec-shape}
For every $\delta_k \in (0, 1)$ and $g_k \in (0, 1)$,
\begin{equation}
  \label{eq:buffer-spec-shape}
  P_{\mathrm{shell}}(\delta_k, g_k)
  \;\le\;
  C_{\mathrm{spec}}(\delta_k) \cdot \exp\!\left(-\frac{\kappa(\SUtwo)}{g_k^2}\right) \cdot g_k^{\beta},
\end{equation}
where
\begin{align}
  \label{eq:spec-shape-data}
  \kappa(\SUtwo) &\;=\; \Casimir(\tfrac12) \;=\; \tfrac34
    \quad \text{(canonical-$\kappa$);}\\
  \beta &\;=\; 3 \;=\; \dim \SUtwo;\\
  C_{\mathrm{spec}}(\delta_k) &\;=\; C_{\mathrm{safety}} \cdot
    \frac{1}{\pi}\int_{\delta_k/2}^{2\,\delta_k} \sin^2(\theta/2)\, d\theta,
\end{align}
with $C_{\mathrm{safety}} = 4$ the empirically-calibrated safety
prefactor that absorbs sub-leading character-tail corrections.
\end{theorem}

\begin{proof}[Proof sketch]
Bound $K_t(\theta)$ on the annulus using the character-tail estimate
of \Cref{ch:heat-kernel}, Proposition~\ref{prop:character-tail},
truncated at $J = \tfrac12$: the leading term is $(2 \cdot \tfrac12 +
1)\,\chi_{1/2}(\theta)\, e^{-t/(3/4)} \le 2\, e^{-t \cdot 3/4}$,
which in the canonical-$\kappa$ identification becomes
$e^{-(3/4)/g_k^2}$.  The Weyl integration of $\sin^2(\theta/2)$ over
the annulus contributes the $\delta_k$-dependent prefactor.  The
$g_k^3$ factor is the shell-volume scaling at the small-$\delta_k$
limit, matching the dimension of $\SUtwo$.  Multiplication by
$C_{\mathrm{safety}}$ absorbs the higher-mode contributions
$(2j+1)^2\,e^{-t\,j(j+1)}$ for $j \ge 1$, which are super-exponentially
suppressed and verified numerically to fit within the safety factor
on the operating range $\delta_k \in (0.02, 1.0)$, $g_k \in (0.05, 1.5)$.
\end{proof}

\marginkey{Uniform-in-$g$ bound}%

\begin{theorem}[Uniform-in-$g$ bound]
\label{thm:buffer-uniform}
For every $\delta_k \in (0.02, 1.0)$ and $g_k \in (0.05, 1.5)$,
\begin{equation}
  \label{eq:buffer-uniform}
  P_{\mathrm{shell}}(\delta_k, g_k)
  \;\le\;
  P_{\mathrm{Haar}}(\delta_k)
  \;+\;
  C_{\mathrm{tail}} \cdot \delta_k^3 \cdot \exp\!\left(-\frac{\kappa(\SUtwo)}{g_k^2}\right),
\end{equation}
where
\begin{align}
  P_{\mathrm{Haar}}(\delta_k) &\;=\;
    \frac{1}{\pi}\int_{\delta_k/2}^{2\,\delta_k} \sin^2(\theta/2)\, d\theta
    \quad\text{(Haar floor)}, \\
  C_{\mathrm{tail}} &\;=\; 4 \quad\text{(uniform-prefactor constant)}.
\end{align}
The bound is uniform: the Haar-floor term controls the large-$g$
limit (where $K_t \to 1$ uniformly), and the
$\delta_k^3\,e^{-(3/4)/g_k^2}$ term controls the small-$g$ tail.
\end{theorem}

\begin{proof}[Proof sketch]
Write $K_t = 1 + (K_t - 1)$ and integrate each summand against the
Weyl measure on the annulus.  The constant-1 contribution gives
$P_{\mathrm{Haar}}(\delta_k)$.  The fluctuation $K_t - 1$ is bounded
above by the irrep-tail of \Cref{ch:heat-kernel},
Proposition~\ref{prop:character-tail}: the $j = \tfrac12$ mode
contributes $2\, e^{-t \cdot 3/4} \cdot |\chi_{1/2}(\theta) - 0| =
2\,|\chi_{1/2}(\theta)|\, e^{-t \cdot 3/4}$, and higher modes are
super-exponentially suppressed.  Integrating
$|\chi_{1/2}(\theta)|\,\sin^2(\theta/2)$ over the annulus contributes
a factor $\Theta(\delta_k^3)$ at small $\delta_k$ (the dominant
small-$\theta$ asymptotic of $|\chi_{1/2}(\theta)| = 2|\cos(\theta/2)|
\approx 2$), and the empirical calibration sets $C_{\mathrm{tail}}
= 4$.  See \cite{Balaban1987RG-I} \S2 for the analogous
$\SUN$ bound; the $\SUtwo$ specialization here picks up the
$\Casimir(\tfrac12) = \tfrac34$ Casimir exponent of
\Cref{ch:su2-rep-theory}.
\end{proof}

\begin{mainfigure}
\centering
\includegraphics[width=\textwidth]{single-plaquette-tail}
\caption{Single-plaquette buffer-shell probability
$P_{\mathrm{shell}}(\delta_k, g_k)$ on $\SUtwo$ at several
$\delta_k$ values across the operating-range $g_k$.  The dashed
curves are the spec-shape bound (\cref{thm:buffer-spec-shape}) and
the uniform-in-$g$ bound (\cref{thm:buffer-uniform}); the solid
markers are the direct heat-kernel quadrature.  The
uniform-in-$g$ bound stays above the quadrature at every
$(\delta_k, g_k)$ in the operating range; the spec-shape bound
over-suppresses the small-$g$ tail through its $g^3$ prefactor.}
\label{fig:single-plaquette-tail}
\end{mainfigure}

\section{Why two bounds}
\label{sec:two-bounds-rationale}

The two bounds serve different regimes and different downstream
consumers.

\begin{itemize}
  \item \textbf{Spec-shape (\cref{thm:buffer-spec-shape}).}  Has the
    explicit $(C, \kappa, \beta)$ form prescribed by the
    overlap-measure aggregate of \Cref{ch:overlap-measure}, where
    the $g_k^3$ factor naturally absorbs into the small-$g$ tail.
    \emph{Not uniform-in-$g$}: at moderate $g_k$ the $g_k^3$ factor
    over-suppresses the tail relative to the actual integral.
  \item \textbf{Uniform-in-$g$ (\cref{thm:buffer-uniform}).}
    Genuinely uniform across the operating range, with the Haar
    floor controlling large-$g$ and the
    $\delta_k^3\,e^{-\kappa/g_k^2}$ controlling small-$g$.  This is
    the bound used by the small-field and overlap KP certificates
    of \Cref{ch:small-field-kp} and \Cref{ch:overlap-kp}.
\end{itemize}

\marginkey{Looseness, motivating chaining}%
Neither bound is optimally tight.  The looseness factor between
the bound and the true heat-kernel integral grows with the
correlation length $\xi(g)$ between adjacent plaquettes; the
per-plaquette union bound on the aggregate overlap measure
(\Cref{ch:overlap-measure}) inherits this looseness multiplicatively.
The chaining-tightening of \Cref{ch:chaining-estimate} addresses
exactly this looseness by exploiting positive correlations among
the per-plaquette events.

\section{Convention dependence}
\label{sec:buffer-convention}

In the Wilson convention of \Cref{ch:heat-kernel},
Convention~\ref{conv:wilson}, the same bounds hold with $\kappa$
replaced by its Wilson value:
\begin{equation}
  \label{eq:wilson-kappa}
  \kappa^{\mathrm{Wilson}}(\SUtwo) \;=\; 2N \cdot \Casimir(\tfrac12) \;=\; 4 \cdot \tfrac34 \;=\; 3 .
\end{equation}
The Manton convention has the wrong $g$-dependence shape and is
\emph{not} accepted by the buffer-shell bound; consumers requiring
a $\exp(-c\, g^2)$ form should instead refer to the diffusion-time
analysis of \cite{Balaban1988RG-II} \S3.

\section{Code provenance}
\label{sec:buffer-codeanchor}

\begin{itemize}
  \item Python implementation:
    \codepath{computations/buffer\_shell.py}.  The module exposes
    \codepath{BufferShellBound} with two methods:
    \codepath{value()} (spec-shape, \cref{thm:buffer-spec-shape})
    and \codepath{uniform\_value()} (uniform-in-$g$,
    \cref{thm:buffer-uniform}); the prefactor constants
    $C_{\mathrm{safety}}$ and $C_{\mathrm{tail}}$ are both set to
    $4$.  Tests in \codepath{tests/test\_buffer\_shell.py}, including
    \codepath{test\_uniform\_value\_is\_uniform\_upper\_bound}
    which verifies \cref{thm:buffer-uniform} on a dense grid in
    $(\delta_k, g_k) \in (0.02, 1.0) \times (0.05, 1.5)$.
  \item Lean formalization: \codepath{lean/YMGap/BufferShell.lean}
    carries the interval-arithmetic-over-$\Qp$ mechanization of the
    bound's $(C, \kappa, \beta)$ shape, with axiom-gate
    \codepath{[propext]} only.
  \item Buffer-shell figure:
    \codepath{visualizations/pilot/su2/single-plaquette-tail.svg}
    shows $P_{\mathrm{shell}}(\delta_k, g_k)$ vs.\ $g_k$ at several
    $\delta_k$ values, with the two bounds overlaid on the direct
    heat-kernel quadrature.
\end{itemize}
